\section{NPE Implementation}\label{ap:npe_implementaion}

We implement Sequential Neural Posterior Estimation (SNPE) using neural spline flows (NSF) as the density estimator, following the framework established by \citet{sbi}. All results presented in this work are fully reproducible using the configuration files and source code available in the project repository at \texttt{https://github.com/tomkimpson/NPE\_for\_RW\_inference}.

For the baseline 1D inference problem, where observations consist of column-wise marginalizations of the spatial distribution, we employ the following architecture and training configuration:
\begin{itemize}
	\item \textbf{Data representation:} Column counts (200 values per observation, corresponding to lattice width $L_x = 200$).
	\item \textbf{Flow architecture:} Neural spline flows (NSF) with 8 coupling layers and 128 hidden features.
	\item \textbf{Training:} Adam optimizer with learning rate $10^{-4}$, batch size 128, training for up to 300 epochs with early stopping based on validation loss.
	\item \textbf{Validation:} 10\% of simulations reserved for validation, with training stopped when validation loss increases for 10 consecutive epochs.
	\item \textbf{SNPE rounds:} 10 rounds with 2,000 simulations per round (10,000 total initial simulations for first round).
	\item \textbf{Convergence:} Sequential rounds terminated when the convergence metric (based on posterior mean and standard deviation changes) falls below 0.001.
	\item \textbf{Implementation:} Built using the SBI library \citep{sbi} with PyTorch backend, trained on NVIDIA A100-SXM4-80GB GPU (85.1 GB memory).
\end{itemize}



\section{CNN implementation} \label{ap:cnn}

To process the full 2D spatial information, we developed a convolutional neural network (CNN) that serves as an embedding network within the NSF architecture. This embedding network compresses the high-dimensional spatial observations into a tractable latent representation that preserves the relevant spatial structure while enabling efficient density estimation. The design incorporates modern architectural elements including residual connections, batch normalization, and progressive downsampling to ensure stable training and effective feature extraction.

The complete CNN architecture comprises the following components:

\begin{itemize}
	\item \textbf{Convolutional layers:} Four progressive downsampling convolutional stages with 32, 64, 128, and 256 filters respectively. Each stage uses $3 \times 3$ kernels with stride 2, batch normalization, ReLU activations, and residual connections implemented via $1 \times 1$ convolutional shortcuts.

	\item \textbf{Dimensionality reduction:} The CNN compresses the $50 \times 200 = 10,000$ dimensional input through progressive downsampling followed by global average pooling, producing a 256-dimensional embedding that captures relevant spatial features while remaining computationally tractable.

	\item \textbf{Fully connected layers:} The pooled convolutional features are passed through three fully connected layers (512, 256, and 256 units) with ReLU activations and dropout (rate = 0.05) for regularization.

	\item \textbf{Normalization:} Per-sample normalization is applied to the input data before processing, with mean subtraction and standard deviation scaling to stabilize training.

	\item \textbf{Translation invariance:} The convolutional architecture with residual connections naturally handles the translation-invariant nature of the spreading process, learning features that are consistent across spatial locations.
\end{itemize}


The integration of the CNN embedding network into the SNPE framework necessitates several important modifications to both the network architecture and training procedure relative to the 1D baseline. These adjustments reflect the substantially increased dimensionality of the input space and the additional complexity introduced by the CNN's learnable parameters. The key differences in our 2D implementation are:

\begin{itemize}
	\item \textbf{Data representation:} Full 2D spatial grid (10,000 values per observation, corresponding to $L_y \times L_x = 50 \times 200$), representing the complete spatial distribution of agents rather than compressed column counts.

	\item \textbf{Flow architecture:} Neural spline flows (NSF) with 5 coupling layers and 256 hidden features (compared to 8 coupling layers and 128 hidden features for 1D data).

	\item \textbf{Increased simulation budget:} Due to the higher-dimensional input space, we use 6,000 simulations per SNPE round (compared to 2,000 for 1D) to ensure adequate coverage of the joint parameter-data space.

	\item \textbf{Extended training time:} The CNN embedding network requires additional optimization iterations. We use a learning rate of $5 \times 10^{-5}$ (compared to $10^{-4}$ for 1D) with batch size 512, training for up to 150 epochs with early stopping after 40 consecutive epochs of no improvement (compared to 10 for 1D).

	\item \textbf{SNPE rounds:} 8 rounds with 6,000 simulations per round (5,000 total initial simulations for first round).

	\item \textbf{Convergence:} Sequential rounds terminated when the convergence metric falls below 0.01 (relaxed from 0.001 for 1D due to the increased complexity of the 2D inference problem).

	\item \textbf{Memory considerations:} Each 2D observation requires significantly more memory ($50 \times 200 = 10,000$ values vs. 200 for column counts), necessitating careful batch size selection and data loading strategies.

	\item \textbf{Implementation:} Same computational infrastructure as 1D (SBI library, PyTorch, NVIDIA A100-SXM4-80GB GPU), with the CNN embedding network integrated directly into the NSF density estimator.
\end{itemize}
